
\svnid{$Id: $}
\svnid{$URL: $}

\section{Junction behaviour for steady state conditions (flow diversion and acceleration at junction; equidistant grid)}
\newrefsegment

%---------------------------------------------------------------------------------
\section*{Quality Assurance}
\begin{tabular}{@{}p{27mm}@{}p{0mm}p{\textwidth-52mm-48pt}p{15mm}p{10mm}}
\textbf{Date} & \textbf{:} & April 26 2019   \\ 
\textbf{Author} & \textbf{:} & Rinske Hutten  & \textbf{Initials:} & RH \\
\textbf{Review} & \textbf{:} &   & \textbf{Initials:} & \ldots 
\end{tabular}

%---------------------------------------------------------------------------------
\section*{Version information}
{{
\begin{tabular}{@{}p{27mm}@{}p{0mm}p{\textwidth-27mm-24pt}}
\textbf{Executable}   & \textbf{:} & Deltares, D-Flow FM Version 1.2.54.64224M, Jun 28 2019, 15:14:31 \\
\textbf{Location input}     & \textbf{:} & \url{https://repos.deltares.nl/repos/DSCTestbench/trunk/cases/e02_dflowfm/f104_1D_numerical-aspects/c01_junction-advection-acceleration-equidistant} \\
\textbf{Revision input} & \textbf{:} & - \\
\textbf{Location output}     & \textbf{:} & \url{https://repos.deltares.nl/repos/DSCTestbench/trunk/references/win64/e02_dflowfm/f104_1D_numerical-aspects/c01_junction-advection-acceleration-equidistant/DFM_OUTPUT_Flow1D} \\
\textbf{Revision output} & \textbf{:} & -
\end{tabular}
}}

%---------------------------------------------------------------------------------
\section*{Purpose}
The purpose of this validation test is to verify that basic network simulations work well.
The test checks that (1) reading a basic network with uniform cross-section and roughness works, (2) simulated discharges and water levels match the imposed constant boundary forcing, (3) a branch split by means of a network (connection) node gives the same result as a single branch, (4) junctions (bifurcations and confluences) conserve mass and are numerically stable, and (5) symmetric networks give symmetric results.

%---------------------------------------------------------------------------------
\section*{Linked claims}
\begin{enumerate}
\item \dflowfm is mass conserving.\footnote{This test case has been copied from \sobekthree test case \textbf{e106}(dflow1d)\_\textbf{f04}(numerical-aspects)\_\textbf{c07}(junction-advection01\_iadvec1d\_1) which was used to validate the implementation of scalar momentum conservation (junction advection).
Since \dflowfm uses a different implementation of which the details are unknown, this test case is currently not yet used to validate the momentum/energy results at the junctions.}
\item \dflowfm can accurately simulate steady and unsteady hydrostatic flow on 1D networks.
\end{enumerate}

%---------------------------------------------------------------------------------
\section*{Approach}
Water levels and velocities are simulated for a number of channel configurations (sub-models); all configurations should give comparable (if not identical) results.
The comparison focuses in particular on the water levels at junctions and velocities near junctions (read: connection nodes).
More precisely, results of sub-models T2 (single channel split by junction), T3 (bifurcating channel) and T4 (combined confluence/bifurcation) are validated against the results of sub-model T1 (without junction).
The results of sub-models T3 and T4, are furthermore compared to three additional reference sub-models to analyze the effect of directional change: T5 (identical to T3, but with the downstream branches on top of each other), T6 (identical to T4, but with the upstream branches on top of each other as well as the downstream ones) and T7 (identical to T4, but with the downstream branch properties flipped).
Finally, the results are also compared to a reference \sobekthree test case.
The sub-models use an equidistant grid and the channel slope increases downstream of the junction location causing an acceleration at the junction.

%---------------------------------------------------------------------------------
\section*{Conclusion}
The network, geometry and forcings are read properly.
The simulated discharges and water levels match the constant forcing at the boundaries.
The results of the model don't change if a branch is split by replacing a regular grid point by a network (connection) node.
Mass is conserved -- also at junctions -- and the inclusion of a junction in the model doesn't cause any instabilities.
The simulation results on a symmetric network are symmetric, but a comparison of the current \dflowfm results to the \sobekthree results for the same geometry shows slight differences in the computed water levels, water depths, discharges and velocities at and around the connection nodes.
The \sobekthree model conserves scalar momentum whereas \dflowfm treats momentum as a vectorial quantity also in a 1D model; as a result changes in flow direction at the junctions influences the results (as shown by an additional set of reference simulations).

%---------------------------------------------------------------------------------
\section*{Model description}
The test comprises of seven sub-models (see \autoref{fig:e02_f104_c01_layout}), all having an \emph{equidistant} grid-spacing of 1000 m.
They all use "open Y-Z rectangular cross-sections" and a constant Ch\'ezy roughness of 50 m\textsuperscript{1/2}/s).
The conveyance is determined using the segmented approach such that the vertical walls don't result in friction.
The longitudinal bed level profile is also the same for all sub-models: an elevation of 10.15 m AD at the upstream boundary/boundaries, 10.00 m AD at 15000 m (the optional confluence/bifurcation point), and 0.00 m AD at 25000 m as depicted in \autoref{fig:e02_f104_c01_profile}.
The downstream boundary is a constant water level of 1.0 m in all cases, and the \emph{total} upstream discharge is always 790.57 m\textsuperscript{3}/s.

\begin{figure}
    \centering
    \includegraphics*[height=0.85\textheight]{figures/e02_f104_c01_layout.jpg}
    \caption{Lay-out of the test model}
    \label{fig:e02_f104_c01_layout}
\end{figure}

\begin{figure}
    \centering
    \includegraphics*[width=0.85\textwidth]{figures/e02_f104_c01_SideView.jpg}
    \caption{Along channel bed profile for all sub-models.
    The junction is always located at the knickpoint (x = 15 km)}
    \label{fig:e02_f104_c01_profile}
\end{figure}

\begin{enumerate}
	\item Sub-model T1 comprises of one branch B1 (length 25000 m) only.
		\begin{itemize}
			\item At T1\_Bnd\_B1\_x0m a constant discharge of 790.57 m\textsuperscript{3}/s is imposed.
			\item The cross-section is uniformly 500 m wide.
			\item The flow accelerates locally at the fixed water-level-point T1\_h\_pnt\_x15000m from 1.13 to 1.56 m/s (see \autoref{fig:e02_f104_c01_profile}).
		\end{itemize}

	\item Sub-model T2 is derived from T1 by splitting the single branch into two parts: an upstream branch B1 (length 15000 m) and a downstream branch B2 (length 10000 m).
	The branches are connected at connection node T2\_ConNode.
		\begin{itemize}
			\item At T2\_Bnd\_B1\_x0m a constant discharge of 790.57 m\textsuperscript{3}/s is imposed.
			\item Both branches have a uniform cross-sectional width of 500 m.
		\end{itemize}

	\item Sub-model T3 is derived from T2 by splitting the downstream branch into two parallel branches of equal length and slope: a right branch B2 and a left branch B3.
	All branches are connected at connection node T3\_ConNode.
		\begin{itemize}
			\item At T3\_Bnd\_B1\_x0m a constant discharge of 790.57 m\textsuperscript{3}/s is imposed.
			\item The upstream branch B1 is still 500 m wide.
			\item The two downstream branches B2 and B3 are both 250 m wide.
		\end{itemize}

		\Note The bed levels and hydraulic properties per meter of cross-sectional width are identical to corresponding sections of sub-model T1.

	\item Sub-model T4 comprises of two upstream branches B1 and B4 (both 15000 m long) and two downstream branches B2 and B3 (both 10000 m long).
	The branches are connected at connection node T4\_ConNode.
		\begin{itemize}
			\item Branch B1 is 150 m wide, and the other upstream branch B4 is 350 m wide (total again 500 m).
			The discharge of 790.57 m\textsuperscript{3}/s is split proportionally over these two branches: 237.171 m\textsuperscript{3}/s enters at T4\_Bnd\_B1\_x0m and 553.399 m\textsuperscript{3}/s at T4\_Bnd\_B4\_x0m.
			\item The downstream branches are no longer symmetric.
			Branch B2 is 50 m wide, and branch B3 is 450 m wide (total again 500 m).
		\end{itemize}

		\Note The main downstream branch B3 (450 m wide) is not aligned with the main upstream branch B4 (350 m wide); compare sub-model T7.

	\item Sub-model T5 is identical to sub-model T3 except that the downstream branches B2 and B3 are now on top of each other and in the same direction as upstream branch \textbf{B1} (aligned with the x-axis).

	\item Sub-model T6 is identical to sub-model T4 except that the upstream branches B1 and B4 are on top of each other and B2 and B3 too.
	The upstream and downstream branches are all aligned with the x-axis.

	\item Sub-model T7 is identical to sub-model T4 except that the widths of the downstream branches B2 and B3 are switched.
	The branch B2 is now the main downstream branch (450 m wide); this branch is aligned with the main upstream branch B4 (350 m wide).
		\begin{itemize}
			\item Branch B1 is 150 m wide, and the other upstream branch B4 is 350 m wide (total again 500 m).
			The discharge of 790.57 m\textsuperscript{3}/s is split proportionally over these two branches: 237.171 m\textsuperscript{3}/s enters at T4\_Bnd\_B1\_x0m and 553.399 m\textsuperscript{3}/s at T4\_Bnd\_B4\_x0m.
			\item The downstream branch branch B2 is 450 m wide, and branch B3 is 50 m wide (total again 500 m).
		\end{itemize}
	\end{enumerate}


%---------------------------------------------------------------------------------
\section*{Results}
Water level, water depth, discharge and velocity results of sub-models T1 to T7 are given in \autoref{tbl:e02_f104_c01_height} and \autoref{tbl:e02_f104_c01_flow} for \dflowfm Version 1.2.48.63929M.
The sub-models T1 to T4 have also been run using \sobekthree (\dflowoned) with junction advection switched on (iadvec1d=1).
The \dflowoned results at the junction are independent of the sub-model: the water depth is always 1.0174 m and hence the water level is always 11.0174.
The discharge and velocity results of \dflowoned are presented in \autoref{tbl:e02_f104_c01_flow}.
Since the \dflowoned results are independent of the channel directions, the \dflowoned results of T3 can also be used for T5, and those of T4 can be used for T6 and T7 (in the latter case with the results for branches B2 and B3 flipped).
Both tables also include columns indicating the difference between \dflowfm and \dflowoned.

\begin{table}[H]
\begin{tabular}{|c|c|cc|c|} \hline
\STRUT & & \multicolumn{2}{c|}{\dflowfm} & $\text{\dflowfm} - \text{\dflowoned}$ \\
\STRUT sub-model & location & water depth \unitbrackets{m} & water level \unitbrackets{m} & difference \unitbrackets{m} \\ [1ex] \hline
\STRUT T1 & T1\_h\_pnt\_x150000m & 1.0117 & 11.0117 & -0.0057 \\
\STRUT T2 & T2\_ConNode          & 1.0117 & 11.0117 & -0.0057 \\
\STRUT T3 & T3\_ConNode          & 1.0144 & 11.0144 & -0.0030 \\
\STRUT T4 & T4\_ConNode          & 1.0224 & 11.0224 &  0.0050 \\
\STRUT T5 & T5\_ConNode          & 1.0129 & 11.0129 & -0.0045 \\
\STRUT T6 & T6\_ConNode          & 1.0120 & 11.0120 & -0.0054 \\
\STRUT T7 & T7\_ConNode          & 1.0177 & 11.0177 &  0.0003 \\ [1ex] \hline
\end{tabular}
    \centering
    \caption{Water levels and water depths of \dflowfm at (connection) nodes.
    Difference compared to \dflowoned.}
    \label{tbl:e02_f104_c01_height}
\end{table}

\begin{table}[H]
\tiny
\begin{tabular}{|c|c|cc|cc|cc|} \hline
\STRUT & & \multicolumn{2}{c|}{\dflowfm} & \multicolumn{2}{c|}{\dflowoned} & \multicolumn{2}{c|}{$\text{\dflowfm} - \text{\dflowoned}$} \\
\STRUT sub-model & location (u-pnt) & discharge \unitbrackets{m\textsuperscript{3}/s} & velocity \unitbrackets{m/s} & discharge \unitbrackets{m\textsuperscript{3}/s} & velocity \unitbrackets{m/s} & difference \unitbrackets{m\textsuperscript{3}/s} & difference \unitbrackets{m/s} \\ [1ex] \hline
\STRUT T1 & B1 before T1\_h\_pnt\_x150000m & 790.570000 & 1.133163 & 790.570000 & 1.133680 &  0.000000 & -0.000517 \\
\STRUT    & B1 after T1\_h\_pnt\_x150000m  & 790.570000 & 1.563655 & 790.570000 & 1.554955 &  0.000000 &  0.008700 \\ [1ex] \hline
\STRUT T2 & B1 next to T2\_ConNode         & 790.570000 & 1.133163 & 790.570000 & 1.133680 &  0.000000 & -0.000517 \\
\STRUT    & B2 next to T2\_ConNode         & 790.570000 & 1.563655 & 790.570000 & 1.554955 &  0.000000 &  0.008700 \\ [1ex] \hline
\STRUT T3 & B1 next to T3\_ConNode         & 790.570000 & 1.136187 & 790.570000 & 1.133680 &  0.000000 &  0.002507 \\
\STRUT    & B2 next to T3\_ConNode         & 395.285000 & 1.559507 & 395.285000 & 1.554955 &  0.000000 &  0.004552 \\
\STRUT    & B3 next to T3\_ConNode         & 395.285000 & 1.559507 & 395.285000 & 1.554955 &  0.000000 &  0.004552 \\ [1ex] \hline
\STRUT T4 & B1 next to T4\_ConNode         & 237.171000 & 1.122178 & 237.171000 & 1.133680 &  0.000000 & -0.011502 \\
\STRUT    & B2 next to T4\_ConNode         &  80.849626 & 1.582450 &  79.057000 & 1.554955 &  1.792626 &  0.027495 \\
\STRUT    & B3 next to T4\_ConNode         & 709.720374 & 1.543465 & 711.513000 & 1.554955 & -1.792626 & -0.011490 \\
\STRUT    & B4 next to T4\_ConNode         & 553.399000 & 1.136441 & 553.399000 & 1.133680 &  0.000000 &  0.002761 \\ [1ex] \hline
\STRUT T5 & B1 next to T5\_ConNode         & 790.570000 & 1.130509 & 790.570000 & 1.133680 &  0.000000 & -0.003171 \\
\STRUT    & B2 next to T5\_ConNode         & 395.285000 & 1.561908 & 395.285000 & 1.554955 &  0.000000 &  0.006953 \\
\STRUT    & B3 next to T5\_ConNode         & 395.285000 & 1.561908 & 395.285000 & 1.554955 &  0.000000 &  0.006953 \\ [1ex] \hline
\STRUT T6 & B1 next to T6\_ConNode         & 237.171000 & 1.128855 & 237.171000 & 1.133680 &  0.000000 & -0.004825 \\
\STRUT    & B2 next to T6\_ConNode         &  78.938846 & 1.560974 &  79.057000 & 1.554955 & -0.118154 &  0.006019 \\
\STRUT    & B3 next to T6\_ConNode         & 711.631154 & 1.563570 & 711.513000 & 1.554955 &  0.118154 &  0.008615 \\
\STRUT    & B4 next to T6\_ConNode         & 553.399000 & 1.131782 & 553.399000 & 1.133680 &  0.000000 & -0.001898 \\ [1ex] \hline
\STRUT T7 & B1 next to T7\_ConNode         & 237.171000 & 1.124221 & 237.171000 & 1.133680 &  0.000000 & -0.009459 \\
\STRUT    & B2 next to T7\_ConNode         & 711.611164 & 1.554696 & 711.513000 & 1.554955 &  0.098164 & -0.000259 \\
\STRUT    & B3 next to T7\_ConNode         &  78.958836 & 1.552551 &  79.057000 & 1.554955 & -0.098164 & -0.002404 \\
\STRUT    & B4 next to T7\_ConNode         & 553.399000 & 1.138797 & 553.399000 & 1.133680 &  0.000000 &  0.005117 \\ [1ex] \hline
\end{tabular}
    \centering
    \caption{Discharges and velocities at u-points directly adjacent to (connection) nodes}
    \label{tbl:e02_f104_c01_flow}
\end{table}

%---------------------------------------------------------------------------------
\section*{Analysis of results}
The results are analysed through (1) an intercomparison of the results of the various sub-models; (2) a comparison of the results with those obtained using a simulation for identical geometry and forcing using \dflowoned using the junction-advection option: \textbf{e106}(dflow1d)\_\textbf{f04}(numerical-aspects)\_\textbf{c07}(junction-advection01\_iadvec1d\_1).
Since the bed at the junction is positioned at 10 m above datum, the water level values are equal to the water depth values plus 10 m: results for water depth and water level are thus consistent.
%The single precision accuracy is 0.000061 for 790.570000.

A comparison of the results of the various sub-models, the following can be observed:

\begin{enumerate}
\item The water depths at the (connection) nodes in sub-models T1 and T2 are identical.
This is logical since the kernel doesn't distinguish between regular computational points and junctions with only two branches connected.
For the same reason, flow-velocities in sub-models T1 and T2, that correspond to locations directly before (or upstream of) and behind (or downstream of) T1\_h\_pnt\_x15000m and T2\_ConNode are identical.
\item The water depth at the junction increases when the change in channel direction is larger (compare sub-model T3 with T5, and T4 and T7 with T6).
This trend is in line with the increase in water depth observed for a single channel system with a directional change at the connection node (see \textbf{e02}(dflowfm)\_\textbf{f104}(numerical-aspects)\_\textbf{c16}(curved-channel-junction-advection-reference).
\item However, even in the perfectly aligned cases (sub-models T5 and T6) the water depths and velocities deviate from the results obtained for a single channel (sub-models T1 and T2), so directional change is not the only aspect of the numerical implementation that plays a role.
\item In the two sub-models T3 and T5 with symmetric bifurcations, the discharge is distributed equally over downstream branches B2 and B3; this indicates that symmetric models result in symmetric results.
However, the both water depths and flow-velocities are influenced by the directional change at the symmetric bifurcation (compare results of sub-models T3 and T5).
\item For the combined confluence and bifurcation (sub-models T4, T6 and T7) the results (water depths, flow-velocities, and discharges) are influenced by both the directional change (compare T7 and T6) and the channel alignment (compare T4 and T7).
The results of the models differ more from the straight single channel results of sub-model T1 as the total change in flow direction increases: zero deflection in T6, main flow directions aligned in T7, main flow direction deflected in T4.
\item For the combined confluence and bifurcation (sub-models T4, T6 and T7) the flow-velocities are influenced by the directional change and channel alignment (see results of sub-models T4 and T7).
The channel alignment also has an influence on the discharge distribution over the downstream branches.
\end{enumerate}

From a comparison with the results of the \dflowoned junction-advection test case \textbf{e106}(dflow1d)\_\textbf{f04}(numerical-aspects)\_\textbf{c07}(junction-advection01\_iadvec1d\_1), the following can be observed:
\begin{enumerate}
\item The water depths at the (connection) nodes are mostly lower for \dflowfm than for \dflowoned.
This difference is the largest for the straight channel.
Other \dflowfm validation test cases show that the results for uniform channels are correct.
\item The downstream discharges are different between \dflowfm and \dflowoned for sub-models T4, T6 and T7.
The \dflowoned results are consistent with the equal distribution of the total discharge (790.57 m\textsuperscript{3}\/s) per unit width.
The \dflowfm results differ from this due to directional aspects of the momentum equation (sub-models T4 and T7) and other numerical aspects (sub-model T6).
\item The differences in the water depths and discharges result also in differences in the computed velocities across branches and sub-models.
\end{enumerate}

Overall, we can summarize that the results of \dflowfm match the are correct at first glance.
Detailed inspection of the results shows some differences in the results compared to \dflowoned that can be attributed to the differences in the implementation of the momentum equation at connection nodes.
However, even for the single channel case (sub-model T1) we observe some differences between \dflowfm and \dflowoned that don't occur in validation cases with a uniform channel; this is likely caused by differences in the momentum discretization as well.
%---------------------------------------------------------------------------------
\printrefsegment

